\documentclass[11pt]{article}
\usepackage[a4paper,margin=1in]{geometry}
\usepackage{amsmath,amssymb,amsthm}
\usepackage[T1]{fontenc}
\usepackage{lmodern}
\usepackage[hidelinks]{hyperref}
\setlength{\emergencystretch}{3em}

\theoremstyle{plain}
\newtheorem{theorem}{Theorem}
\newtheorem{lemma}[theorem]{Lemma}
\newtheorem{proposition}[theorem]{Proposition}
\newtheorem{corollary}[theorem]{Corollary}
\theoremstyle{definition}
\newtheorem{definition}[theorem]{Definition}
\theoremstyle{remark}
\newtheorem{remark}[theorem]{Remark}

\title{A Proof of Conjecture 00000000051:\\
the Graphs $K_{p,p}$ Have Spectrum $\{0,\pm p\}$}
\author{%
  Feiyu\thanks{Independent researcher.}
  \and
  Claude (Opus 5.5)\thanks{Anthropic.}%
}
\date{2026-10-02}

\begin{document}
\maketitle

\begin{abstract}
Conjecture 00000000051 asserts that there are infinitely many pairwise
non-isomorphic connected integral graphs whose adjacency spectrum consists only
of $0$ and numbers $\pm p$ with $p$ prime. The complete bipartite graphs
$K_{p,p}$, $p$ prime, form such a family: their adjacency matrix $A$ satisfies
$A^3=p^2A$, so every eigenvalue lies in $\{0,p,-p\}$. The proof is formalized
in Lean~4 against Mathlib.
\end{abstract}

\section{The conjecture as stated}

The original text of conjecture 00000000051 reads:

\begin{quote}
Conjecture: There exist infinitely many pairwise non-isomorphic connected
integral graphs whose adjacency spectrum consists only of $0$ or $\pm$primes.
\end{quote}

\section{Reading of the statement}

Graphs are finite and simple. The \emph{adjacency spectrum} of a graph $G$ is
the set of (complex) eigenvalues of its adjacency matrix $A(G)$. As usual, $G$
is \emph{integral} if all its adjacency eigenvalues are integers.

\begin{definition}
A finite simple graph $G$ is \emph{good} if it is connected, integral, and
every adjacency eigenvalue $\mu$ of $G$ satisfies $\mu=0$ or $\mu=\pm p$ for
some prime $p$.
\end{definition}

The conjecture asserts that there are infinitely many isomorphism classes of
good graphs, that is, an infinite sequence of good graphs no two of which are
isomorphic. (Integrality is implied by the spectral condition, since $0$ and
$\pm p$ are integers, but we keep it as stated.)

\section{Result}

\begin{theorem}\label{thm:main}
For every prime $p$ the complete bipartite graph $K_{p,p}$ is good. Since
$K_{p,p}$ has $2p$ vertices, these graphs are pairwise non-isomorphic for
distinct primes, and as there are infinitely many primes, Conjecture
00000000051 holds.
\end{theorem}

\section{Proof}

\begin{lemma}\label{lem:cube}
Let $A$ be an $N\times N$ complex matrix and $c\in\mathbb C$ with $A^3=c^2A$.
Then every eigenvalue of $A$ lies in $\{0,c,-c\}$.
\end{lemma}

\begin{proof}
If $Av=\mu v$ with $v\ne0$, then $\mu^3v=A^3v=c^2Av=c^2\mu v$, so
$\mu^3-c^2\mu=\mu(\mu-c)(\mu+c)=0$.
\end{proof}

\begin{lemma}\label{lem:walks}
Let $A$ be the adjacency matrix of $K_{m,m}$ with sides $L$ and $R$. Then
$(A^2)_{ij}=m$ if $i,j$ lie on the same side and $(A^2)_{ij}=0$ otherwise.
Consequently $A^3=m^2A$.
\end{lemma}

\begin{proof}
$(A^2)_{ij}$ counts vertices $k$ adjacent to both $i$ and $j$, that is, lying on
the side opposite to both. There are $m$ of them if $i$ and $j$ are on the same
side, and none otherwise. Then
$(A^3)_{ij}=\sum_k A_{ik}(A^2)_{kj}$ runs over the $m$ vertices $k$ opposite to
$i$, each contributing $m$ exactly when $k$ is on the side of $j$, i.e.\ when
$j$ is opposite to $i$. So $(A^3)_{ij}=m^2A_{ij}$.
\end{proof}

\begin{proof}[Proof of Theorem~\ref{thm:main}]
Let $p$ be prime and $A$ the adjacency matrix of $K_{p,p}$. By
Lemmas~\ref{lem:walks} and~\ref{lem:cube}, every eigenvalue of $A$ is $0$, $p$
or $-p$. These are integers, so $K_{p,p}$ is integral, and they are of the form
required by the conjecture. $K_{p,p}$ is connected, since $p\ge1$: two vertices
on opposite sides are adjacent, and two vertices on the same side have a common
neighbour. Finally, an isomorphism $K_{p,p}\cong K_{q,q}$ is a bijection
between vertex sets of sizes $2p$ and $2q$, so $p=q$. Enumerating the primes as
$p_0<p_1<\cdots$ gives the infinite sequence $K_{p_0,p_0},K_{p_1,p_1},\ldots$
of pairwise non-isomorphic good graphs.
\end{proof}

\begin{remark}
In fact the spectrum of $K_{p,p}$ is exactly $\{p,-p,0\}$, with $0$ of
multiplicity $2p-2$: the vectors that are constant on each side span an
$A$-invariant plane on which $A$ acts by
$\left(\begin{smallmatrix}0&p\\p&0\end{smallmatrix}\right)$. Only the inclusion
is needed for the conjecture.
\end{remark}

\section{Formalization}

The accompanying Lean~4 project (\texttt{lean/Submission/Basic.lean}) states
the conjecture and proves it against Mathlib.

\begin{itemize}
  \item \texttt{adjSpectrum G} --- Mathlib's \texttt{spectrum} of the adjacency
    matrix \texttt{G.adjMatrix} over $\mathbb C$, taken in the matrix algebra.
    For a matrix this is exactly its set of eigenvalues.
  \item \texttt{IsIntegralGraph}, \texttt{SpectrumZeroOrPmPrime},
    \texttt{IsGood} --- the definitions above, with Mathlib's
    \texttt{SimpleGraph.Connected} and \texttt{Nat.Prime}.
  \item \texttt{ConjectureHolds} --- there is a sequence of finite graphs
    $(G_k)_{k\in\mathbb N}$, each \texttt{IsGood}, such that the type of graph
    isomorphisms $G_j\simeq G_k$ is empty whenever $j\ne k$.
  \item \texttt{spectrum\_subset\_of\_cube} --- Lemma~\ref{lem:cube}, via
    Mathlib's spectral mapping inclusion
    \texttt{spectrum.subset\_polynomial\_aeval} applied to $X^3-c^2X$.
  \item \texttt{adjMatrix\_sq\_apply}, \texttt{adjMatrix\_cube} ---
    Lemma~\ref{lem:walks}, for $K_{m,m}$ taken as Mathlib's
    \texttt{completeBipartiteGraph} on two copies of \texttt{Fin m}.
  \item \texttt{adjSpectrum\_K}, \texttt{K\_connected}, \texttt{isGood\_K} ---
    $K_{p,p}$ is good.
  \item \texttt{conjecture\_00000000051} --- \texttt{ConjectureHolds}, with
    $G_k=K_{p_k,p_k}$ for $p_k=\texttt{Nat.nth Nat.Prime k}$; distinct indices
    give distinct primes (\texttt{Nat.nth\_injective}), hence different vertex
    counts.
\end{itemize}

The toolchain is \texttt{leanprover/lean4:v4.33.1}, Mathlib is pinned to
\texttt{v4.33.1}, and \texttt{lake build} succeeds. The project contains no
\texttt{sorry}, no \texttt{admit} and no \texttt{native\_decide}, and
introduces no axioms:
\begin{center}
\texttt{\#print axioms conjecture\_00000000051}
\end{center}
reports exactly \texttt{propext}, \texttt{Classical.choice} and
\texttt{Quot.sound}.

\end{document}
